
\section*{\underline{EXPOSÉ VIII}\\[3pt] \underline{\underline{LE THÉORÈME DE RIEMANN--ROCH}}}
\begin{center}
\underline{par Pierre BERTHELOT}
\end{center}

\subsection*{1. Morphismes d'intersection complète}

\textbf{Définition 1.1.} \uline{Soit} $f:X\to Y$ \uline{un morphisme de schémas}. \uline{On dit que} $f$ \uline{est un morphisme d'intersection complète si pour tout point} $x$ \uline{de} $X$, \uline{il existe un voisinage ouvert} $U$ \uline{de} $x$, \uline{un} $Y$-\uline{schéma lisse} $V$, \uline{tels que la restriction de} $f$ \uline{à} $U$ \uline{se factorise en} :
\[
\begin{tikzcd}[column sep=large,row sep=large]
U \arrow[rr,"i"] \arrow[dr] && V \arrow[dl]\\
& Y &
\end{tikzcd}
\]
\uline{où} $i$ \uline{est une immersion régulière} (VII 1.4).

Quitte à remplacer $V$ par un ouvert, on peut supposer que $i$ est une immersion fermée régulière.

On peut dans cette définition remplacer $V$ par un fibré vectoriel $Y[T_1,\ldots,T_n]$. En effet, la notion étant locale, on peut supposer que $V$ est affine, et que son image dans $Y$ est contenue dans un ouvert affine $W$ de $Y$. Il existe alors un entier $n$ et une $W$-immersion fermée $V\longrightarrow W[T_1,\ldots,T_n]$, d'où une $Y$-immersion $V\to Y[T_1,\ldots,T_n]$, qui est régulière d'après VII 1.10. Comme la composée de deux immersions régulières est régulière, on trouve donc une $Y$-immersion régulière $U\longrightarrow Y[T_1,\ldots,T_n]$ factorisant $f$.

\textbf{Proposition 1.2.} \uline{Soit} $f:X\to Y$ \uline{un morphisme d'intersection complète}. \uline{Alors toute} $Y$-\uline{immersion} $i:X\to Z$ \uline{de} $X$ \uline{dans un} $Y$-\uline{schéma lisse est régulière}.

Nous montrerons d'abord le corollaire :

\textbf{Corollaire 1.3.} \uline{Soit le diagramme commutatif} :
\[
\begin{tikzcd}[column sep=large,row sep=large]
& V' \arrow[d,"p"]\\
X \arrow[ur,"j"] \arrow[r,"i"] & V ,
\end{tikzcd}
\]
